% Composición focal ES/EN. Antecedentes conservados y pruebas contiguas.
\section{\HMTLang{Retour longitudinal, réciprocité radiale et archive de mémoire}{Longitudinal return, radial reciprocity, and the memory archive}}
\label{xii:radial:section}

\HMTLang{La constante de couplage conserve son registre ordonné et ses transformations inversibles. La réalisation suivante compose ces opérateurs avec les incidences déjà construites, la norme locale d’action et la réponse constitutive. L’état APP--TRIT--TPK et la structure discrète conjointe du continu précèdent leurs deux lectures. Les coordonnées \(\pi,\varphi,e,\alpha\) reçoivent la généalogie des sections précédentes ; le rayon dimensionnel est construit ensuite sur une section de longueur explicite.}{The coupling constant retains its ordered register and invertible transformations. The following realization composes those operators with the incidences already constructed, the local action norm, and the constitutive response. The APP--TRIT--TPK state and the joint discrete structure of the continuum precede their two readouts. The coordinates \(\pi,\varphi,e,\alpha\) inherit the genealogy established in the preceding sections; the dimensional radius is subsequently constructed on an explicit length section.}

\subsection{\HMTLang{Restitution constitutive et moment orienté}{Restitution constitutive et moment orienté}}
\HMTLang{La réponse de~\eqref{iii:eq:rchannels} et la fonctionnelle de~\eqref{v:ant:eq:gamma} partagent les canaux étiquetés. Pour \(s=q^{30}\), elles conservent la fonction}{La réponse dans~\eqref{iii:eq:rchannels} et la fonctionnelle dans~\eqref{v:ant:eq:gamma} partagent les canaux étiquetés. Pour \(s=q^{30}\), elles conservent la fonction}
\[
 f(s)=\frac{1+s+s^2}{1+s+s^2+s^3},\qquad 0<s<1.
\]
\HMTLang{Le moment orienté de profondeur deux et son opérateur différentiel sont}{Le moment orienté de profondeur deux et son opérateur différentiel sont}

\[
 \mathcal C_2(s)=\sum_{k\ge0}\frac{(-1)^ks^{2k+1}}{(2k+1)^2},
 \qquad \mathcal D=s\frac{d}{ds}.
\]
\HMTLang{La convergence locale uniforme donne \(\mathcal D^2\mathcal C_2=s/(1+s^2)\), et donc}{La convergence locale uniforme donne \(\mathcal D^2\mathcal C_2=s/(1+s^2)\), d’où}
\begin{equation}
 f(s)=\frac{1+s+s^2}{s(1+s)}\mathcal D^2\mathcal C_2(s),
 \qquad r_\pm=f(q_\pm^{30}).
 \label{xii:radial:catalan}
\end{equation}
\HMTLang{La condition \(\mathcal C_2(s)\to0\) lorsque \(s\downarrow0\) fixe la restitution}{La condition \(\mathcal C_2(s)\to0\) lorsque \(s\downarrow0\) fixe la restitution}
\[
 \mathcal C_2(s)=\int_0^s\log(s/t)\frac{1+t}{1+t+t^2}f(t)\,dt.
\]
\HMTLang{En effet, intégrer deux fois par rapport à \(\log s\) inverse \(\mathcal D^2\) ; la différence de deux solutions est \(a\log s+b\), et la limite nulle impose \(a=b=0\). L’intégrale se réduit à \(\int_0^s\log(s/t)/(1+t^2)\,dt\). La série originale converge absolument en \(s=1\) et fournit la constante de Catalan ; sa dérivée seconde à l’extrémité s’interprète par la limite d’Abel. On conserve ainsi la fonction paramétrique complète, en plus de son évaluation scalaire.}{En effet, intégrer deux fois par rapport à \(\log s\) inverse \(\mathcal D^2\) ; la différence de deux solutions est \(a\log s+b\), et la limite nulle impose \(a=b=0\). L’intégrale se réduit à \(\int_0^s\log(s/t)/(1+t^2)\,dt\). La série originale converge absolument en \(s=1\) et donne la constante de Catalan ; sa dérivée seconde à l’extrémité s’interprète par la limite d’Abel. La fonction paramétrique complète est ainsi conservée avec son évaluation scalaire.}


\HMTLang{Le système de coordonnées analytique conserve les représentants canoniques de la chambre fondamentale : \(x=\pi A_{\rm deg}/180\), \(y=\pi C^*_{\rm deg}/180\) et \(q_\pm=e^{-x\mp y}\). Ces représentants relevés fixent les exponentielles réelles. Avec \(\mathcal D_2(z)=\sum_{n\ge1}z^n/n^2\), le potentiel commun est}{The analytic chart retains the canonical representatives of the fundamental chamber: \(x=\pi A_{\rm deg}/180\), \(y=\pi C^*_{\rm deg}/180\), and \(q_\pm=e^{-x\mp y}\). These lifted representatives fix the real exponentials. With \(\mathcal D_2(z)=\sum_{n\ge1}z^n/n^2\), the common potential is}
\[
 \mathscr F(x,y)=\sum_{\sigma=\pm}
 \left[\frac{\mathcal D_2(q_\sigma^{120})}{120}
 -\frac{\mathcal D_2(q_\sigma^{90})}{90}\right],\qquad x>|y|.
\]
\HMTLang{L’identité \(z\mathcal D_2'(z)=-\log(1-z)\) prouve \(\partial_x\mathscr F=\log\widehat c\) et \(\partial_y\mathscr F=\log\widehat Z\). Sa hessienne symétrique conserve \(\partial_y\log\widehat c=\partial_x\log\widehat Z\) ; la condition \(\mathscr F\to0\) lorsque \(q_+,q_-\to0\) fixe la constante additive.}{L’identité \(z\mathcal D_2'(z)=-\log(1-z)\) prouve \(\partial_x\mathscr F=\log\widehat c\) et \(\partial_y\mathscr F=\log\widehat Z\). Sa hessienne symétrique conserve \(\partial_y\log\widehat c=\partial_x\log\widehat Z\) ; la condition \(\mathscr F\to0\) lorsque \(q_+,q_-\to0\) fixe la constante additive.}

\begin{theorem}[\HMTLang{Coordonnées constitutives conjointes}{Joint constitutive coordinates}]
\label{xii:radial:constitutive-inverse}
\HMTLang{L’application \((r_+,r_-)\mapsto(\widehat c,\gamma)\) est un difféomorphisme de \((3/4,1)^2\) sur \((1,16/9)\times\mathbb R\). Elle conserve l’attribution des feuillets.}{The map \((r_+,r_-)\mapsto(\widehat c,\gamma)\) is a diffeomorphism from \((3/4,1)^2\) onto \((1,16/9)\times\mathbb R\). It preserves the assignment of sheets.}
\end{theorem}
\begin{proof}
\HMTLang{La dérivée de \(f\) dans~\eqref{iii:eq:monotone} est strictement négative ; soit \(\psi=f^{-1}\). Définissons}{The derivative of \(f\) in~\eqref{iii:eq:monotone} is strictly negative; let \(\psi=f^{-1}\). Define}
\[
 \mathcal H_B(s)=\frac{12s^3}{1-s^9}-\frac{s^4}{1-s^{12}}
 -\frac{(\log s)^2}{20\pi},\qquad F_B=\mathcal H_B\circ\psi.
\]
\HMTLang{On vérifie \(\gamma=F_B(r_-)-F_B(r_+)\). De plus,}{On vérifie \(\gamma=F_B(r_-)-F_B(r_+)\). De plus,}
\[
 \mathcal H_B'(s)=\frac{36s^2(1+2s^9)}{(1-s^9)^2}
 -\frac{4s^3(1+2s^{12})}{(1-s^{12})^2}-\frac{\log s}{10\pi s}>0.
\]
\HMTLang{Le quotient des deux premiers termes positifs, avant la soustraction, est supérieur à \(9\) ; le dernier terme de la somme est positif. Par conséquent, \(F_B\) décroît de \(+\infty\) à \(-\infty\). \(v=\widehat c\) étant fixé, écrivons \(r_\pm=v^{-1/2}e^{\mp z/2}\). L’intervalle admissible est \(|z|<\min\{\log v,\log(16/(9v))\}\). La dérivée de \(F_B(r_-)-F_B(r_+)\) est \([r_-F_B'(r_-)+r_+F_B'(r_+)]/2<0\) ; ses limites sont \(+\infty\) et \(-\infty\). Il existe une solution unique pour chaque \(\gamma\), avec un inverse lisse grâce à cette dérivée non nulle.}{Le rapport des deux premiers termes positifs, avant la soustraction, dépasse \(9\) ; le dernier terme de la somme est positif. Ainsi, \(F_B\) décroît de \(+\infty\) à \(-\infty\). Pour \(v=\widehat c\) fixé, écrivons \(r_\pm=v^{-1/2}e^{\mp z/2}\). L’intervalle admissible est \(|z|<\min\{\log v,\log(16/(9v))\}\). La dérivée de \(F_B(r_-)-F_B(r_+)\) est \([r_-F_B'(r_-)+r_+F_B'(r_+)]/2<0\), et ses limites sont \(+\infty\) et \(-\infty\). Il existe une solution unique pour chaque \(\gamma\), avec un inverse lisse parce que cette dérivée est non nulle.}
\end{proof}
\HMTLang{Sur l’image engendrée, \(s_\pm=\psi(r_\pm)\) récupère les coordonnées sexagésimales \(A_{\rm deg}=-3\log(s_+s_-)/\pi\), \(C^*_{\rm deg}=3\log(s_-/s_+)/\pi\), puis \(\alpha=A_{\rm deg}/1000\), \(\eta_{\rm ret}=C^*_{\rm deg}/2-\alpha\). Cette dernière opération est l’inverse de la section d’action déjà construite. La récupération fixe les canaux ; les marques, les unités et le lecteur longitudinal conservent leurs fonctions propres.}{On the generated image, \(s_\pm=\psi(r_\pm)\) recovers the sexagesimal coordinates \(A_{\rm deg}=-3\log(s_+s_-)/\pi\), \(C^*_{\rm deg}=3\log(s_-/s_+)/\pi\), followed by \(\alpha=A_{\rm deg}/1000\) and \(\eta_{\rm ret}=C^*_{\rm deg}/2-\alpha\). The last operation inverts the action section already constructed. This recovery fixes the channels; the markings, units, and longitudinal reader retain their own roles.}

\subsection{\HMTLang{Inscription marquée et retour linéaire}{Inscription marquée et retour linéaire}
}
\HMTLang{L’incidence hexadique--octadique et les orbites dodécaphasiques définissent le support}{L’incidence hexadique--octadique et les orbites dodécaphasiques définissent le support}

\[
 \Omega=\{(j,k,q,u,v):j,k\in\mathbb Z/12\mathbb Z,\quad
 k-j\in\{0,3,6,9\},\quad 1\le q\le8,
 \quad u<v,\quad u,v\in\{1,2,4,5,7,8\}\}.
\]
\HMTLang{Son cardinal est \(N=12\cdot4\cdot8\binom62=5760\). La différence \(k-j\), calculée dans \(\mathbb Z/12\mathbb Z\), l’indice d’octade et la paire d’hexade constituent les \(480\) étiquettes auxiliaires par secteur. Sur cet ordre, l’inverseur entier de la section~\ref{subsec:k-hadamard} se prolonge au moyen de}{Son cardinal est \(N=12\cdot4\cdot8\binom62=5760\). La différence \(k-j\), calculée dans \(\mathbb Z/12\mathbb Z\), l’indice d’octade et la paire d’hexade fournissent les \(480\) étiquettes auxiliaires par secteur. Dans cet ordre, l’inverse entier de la section~\ref{subsec:k-hadamard} se prolonge sous la forme}

\[
 B_\Omega=(\mathcal H_{12}/4)\otimes I_{480},\qquad
 B_\Omega^*B_\Omega=B_\Omega B_\Omega^*=\tfrac14 I.
\]
\HMTLang{Les identités \(\mathcal H_{12}^*=\mathcal H_{12}\) et \(\mathcal H_{12}^2=4I\), obtenues à partir des trois blocs de Hadamard d’ordre quatre et de leur permutation conjointe, prouvent l’égalité et conservent l’orientation des trois orbites. Soit \(u=N^{-1/2}\mathbf1\), \(P=uu^*\), \(C_0=(I-P)B_\Omega\), \(M_0=PB_\Omega\) et \(U=2B_\Omega\). Alors}{The identities \(\mathcal H_{12}^*=\mathcal H_{12}\) and \(\mathcal H_{12}^2=4I\), obtained from the three order-four Hadamard blocks and their joint permutation, prove the equality and preserve the orientation of the three orbits. Set \(u=N^{-1/2}\mathbf1\), \(P=uu^*\), \(C_0=(I-P)B_\Omega\), \(M_0=PB_\Omega\), and \(U=2B_\Omega\). Then}
\[
 C_0^*C_0+M_0^*M_0=\tfrac14 I,\qquad
 C_0C_0^*=\tfrac14(I-P),\qquad U^*U=UU^*=I.
\]
\HMTLang{L’archive \((2C_0v,2M_0v)\) est isométrique et sa somme récupère \(Uv\). Le transport \(U\) appartient à ce registre fini ; l’évolution \(U_t\) de TPK conserve en outre la profondeur et l’histoire de la section.}{L’archive \((2C_0v,2M_0v)\) est isométrique, et sa somme récupère \(Uv\). Le transport \(U\) agit sur ce registre fini ; l’évolution TPK \(U_t\) conserve en outre la profondeur et l’histoire de la section.}


\begin{proposition}[\HMTLang{Réponse de l’inscription marquée}{Réponse de l’inscription marquée}]
\label{xii:radial:diagonal}
\HMTLang{Avec les projecteurs atomiques \(P_i=e_ie_i^*\), l’unique espérance sur leur algèbre diagonale qui fixe chaque \(P_i\) et est bimodulaire est}{Pour les projections atomiques \(P_i=e_ie_i^*\), l’unique espérance sur leur algèbre diagonale qui fixe chaque \(P_i\) et est bimodulaire est}
\[
 \Delta_\Omega(T)=\sum_iP_iTP_i,
 \qquad \Delta_\Omega(C_0C_0^*)=r_NI,
 \qquad r_N=\frac{N-1}{4N}=\frac{5759}{23040}.
\]
\end{proposition}
\begin{proof}
\HMTLang{Sur chaque unité matricielle \(E_{ij}\), la bimodularité impose \(\Delta(E_{ij})=P_i\Delta(E_{ij})P_j\). La diagonale annule cette expression si \(i\ne j\) ; fixer \(P_i\) détermine le cas restant. La linéarité fixe toute l’application. Puisque \(P_{ii}=1/N\), la diagonale de \((I-P)/4\) vaut \(r_NI\). La composante commune apporte \(I/(4N)\) ; leur somme est \(I/4\).}{On each matrix unit \(E_{ij}\), bimodularity imposes \(\Delta(E_{ij})=P_i\Delta(E_{ij})P_j\). Diagonality makes this expression zero when \(i\ne j\); fixing \(P_i\) determines the remaining case. Linearity fixes the whole map. Since \(P_{ii}=1/N\), the diagonal of \((I-P)/4\) is \(r_NI\). The common component contributes \(I/(4N)\); the two add to \(I/4\).}
\end{proof}
\HMTLang{L’isométrie \(We_i=e_i\otimes e_i\) réalise \(W^*(T\otimes I)W=\Delta_\Omega(T)\) en conservant le registre élargi. L’uniformité concerne ici le mode commun ; tout état sur l’algèbre marquée évalue le même scalaire \(r_N\). Une mémoire auxiliaire se transporte au moyen de \(\Delta_\Omega\otimes\operatorname{id}\).}{L’isométrie \(We_i=e_i\otimes e_i\) réalise \(W^*(T\otimes I)W=\Delta_\Omega(T)\) en conservant le registre étendu. L’uniformité concerne ici le mode commun ; tout état sur l’algèbre marquée évalue le même scalaire \(r_N\). Une mémoire auxiliaire se transporte au moyen de \(\Delta_\Omega\otimes\operatorname{id}\).}


\begin{definition}[\HMTLang{Composition longitudinale du retour}{Composition longitudinale du retour}]
\label{xii:radial:longitudinal}
\HMTLang{Sur une référence \(L_*>0\) de la droite de longueur, la norme locale \(\alpha^{16}\), construite au moyen du conoyau de Hadamard dans la section d’action, agit linéairement sur le retour inscrit :}{Sur une référence \(L_*>0\) dans la droite de longueur, la norme locale \(\alpha^{16}\), construite au moyen du conoyau de Hadamard dans la section d’action, agit linéairement sur le retour inscrit :}
\begin{equation}
 D=L_*\alpha^{16}\Delta_\Omega(C_0C_0^*)U=LU,
 \qquad L=L_*\alpha^{16}r_N.
 \label{xii:radial:length}
\end{equation}
\end{definition}
\HMTLang{Cette composition réunit l’inverseur, son archive complémentaire et la section longitudinale en une règle explicite. Pour une cochaîne orientée \(\beta\), elle multiplie par \(\alpha^{16}r_N\) ; elle conserve la somme, l’inversion d’orientation et la subdivision en neuf arêtes. La norme de \(C_0^*ae_i\) est \(|a|\sqrt{r_N}\), tandis que la réponse linéaire inscrite est \(ar_N\) : ce sont des opérations différentes. Le coefficient de~\eqref{xii:radial:length} provient de l’espérance spécifiée, outre son type linéaire.}{Cette composition réunit l’inverse, son archive complémentaire et la section longitudinale en une règle explicite. Sur une cochaîne orientée \(\beta\), elle multiplie par \(\alpha^{16}r_N\), en conservant l’addition, l’inversion d’orientation et la subdivision en neuf arêtes. La norme de \(C_0^*ae_i\) est \(|a|\sqrt{r_N}\), tandis que la réponse linéaire inscrite est \(ar_N\) : ce sont des opérations différentes. Le coefficient dans~\eqref{xii:radial:length} découle de l’espérance spécifiée ainsi que de son type linéaire.}

\subsection{\HMTLang{Module d’arrivée et couplage commun}{Module d’arrivée et couplage commun}
}
\HMTLang{Soient \(X=X^*>0\) le caractère complet dans la fibre finie non nulle du registre et \(Y=UXU^*\). La lecture radiale de cette réalisation est le module d’arrivée}{Soient \(X=X^*>0\) le caractère complet sur la fibre finie non nulle du registre et \(Y=UXU^*\). La lecture radiale de cette réalisation est le module d’arrivée}

\[
 R_X=(DX^2D^*)^{1/2}=LY.
\]
\HMTLang{Le radion conserve la loi d’action \(S(\theta)=\hbar_a\theta\). Dans la réalisation circulaire du même \(L\), \(\omega L=c\), de sorte que}{The radion retains the action law \(S(\theta)=\hbar_a\theta\). In the circular realization of the same \(L\), \(\omega L=c\), hence}
\begin{equation}
 Q_a=\frac{\hbar_a c}{L},\quad
 H_X=Q_aY,\quad M_X=H_X/c^2,\quad \Lambda_X=LY^{-1}.
 \label{xii:radial:readouts}
\end{equation}
\HMTLang{La lettre \(Q_a\) désigne ici une énergie ; les projecteurs et les registres précédents conservent leur notation propre. L’identification physique de \(R_X\) est celle du rayon gravitationnel réduit \(R_X=(G_a/c^2)M_X\). Cette identification et la composition longitudinale définissent la réalisation dont les coefficients sont déterminés ci-après.}{Here \(Q_a\) denotes an energy; the preceding projections and registers retain their own notation. The physical identification of \(R_X\) is the reduced gravitational radius \(R_X=(G_a/c^2)M_X\). This identification and the longitudinal composition define the realization whose coefficients are determined below.}

\begin{theorem}[\HMTLang{Rigidité radiale et réciprocité normalisée}{Rigidité radiale et réciprocité normalisée}
]
\label{xii:radial:rigidity}
\HMTLang{Pour \(D\) et \(X\) fixés, \(R_X\) est l’unique opérateur positif qui satisfait \(\|R_Xv\|^2=\|XD^*v\|^2\) pour tout vecteur d’arrivée. Les lectures précédentes satisfont}{Pour \(D\) et \(X\) fixés, \(R_X\) est l’unique opérateur positif satisfaisant \(\|R_Xv\|^2=\|XD^*v\|^2\) pour tout vecteur d’arrivée. Les lectures précédentes satisfont}

\begin{equation}
 R_X\Lambda_X=L^2I,\qquad H_X\Lambda_X=\hbar_a cI,
 \qquad \boxed{G_a=\frac{c^3L^2}{\hbar_a}}.
 \label{xii:radial:G}
\end{equation}
\HMTLang{Le même \(G_a\) agit sur tous les caractères positifs admissibles de cette réalisation.}{Le même \(G_a\) agit sur tout caractère positif admissible de cette réalisation.}

\end{theorem}
\begin{proof}
\HMTLang{La polarisation de l’égalité des formes donne \(R_X^2=DX^2D^*=L^2UX^2U^*\). L’unique racine positive est \(LUXU^*\). Les produits de~\eqref{xii:radial:readouts} donnent les deux identités. Les multiplier par \(\Lambda_X^{-1}\) à droite prouve \(R_X=L^2H_X/(\hbar_a c)\) ; cette déduction conserve tous les termes croisés. Substituer \(H_X=c^2M_X\) et l’identification radiale produit~\eqref{xii:radial:G}. L’inversibilité de \(M_X\) fixe le coefficient.}{La polarisation de l’identité des formes quadratiques donne \(R_X^2=DX^2D^*=L^2UX^2U^*\). Son unique racine carrée positive est \(LUXU^*\). Les produits dans~\eqref{xii:radial:readouts} donnent les deux identités. La multiplication à droite par \(\Lambda_X^{-1}\) prouve \(R_X=L^2H_X/(\hbar_a c)\) ; cette déduction conserve chaque terme croisé. Substituer \(H_X=c^2M_X\) et l’identification radiale donne~\eqref{xii:radial:G}. L’inversibilité de \(M_X\) fixe le coefficient.}
\end{proof}
\HMTLang{L’unicité de l’opérateur correspond à chaque \(X\) fixé, tandis que celle du coefficient est commune à toute la fibre. Le secteur nul conserve son traitement séparé de l’inverse. La loi native \(L_{-p}(z^{-1})=-L_p(z)\), avec \(L_p(z)=(z^p-1)/p\) et \(L_0(z)=\log z\), exprime l’inversion orientée avant cette réalisation gravitationnelle ; la preuve est la substitution directe de \((z^{-1})^{-p}=z^p\).}{Operator uniqueness holds for each fixed \(X\), whereas the coefficient is common to the whole fibre. The null sector retains its separate treatment outside the inverse operation. The native law \(L_{-p}(z^{-1})=-L_p(z)\), with \(L_p(z)=(z^p-1)/p\) and \(L_0(z)=\log z\), expresses oriented inversion before this gravitational realization; its proof is the direct substitution \((z^{-1})^{-p}=z^p\).}

\subsection{\HMTLang{Horloge, masse et sections dimensionnelles}{Horloge, masse et sections dimensionnelles}
}
\HMTLang{Le système de coordonnées chronologique compatible satisfait}{The compatible chronological chart satisfies}
\begin{equation}
 108\ell_0=2\pi L,\quad \ell_0=ct_0=\widehat c\,u_L,
 \quad t_0=\frac{\pi L}{54c},\quad
 \frac{L_*}{u_L}=\frac{54\widehat c}{\pi\alpha^{16}r_N}.
 \label{xii:radial:clock}
\end{equation}
\HMTLang{On distingue ainsi la référence radiale \(L_*\), la base \(u_L=u_Cu_T\) et l’avancée \(\ell_0\). La période \(108t_0\) est \(2\pi L/c\). Le mode \(X=I\) détermine le quantum de l’horloge ; la coordonnée électronique conserve en outre sa normalisation \(b_e=E_e^{(0)}/Q_{\rm ret}\) et son facteur \(R_{\rm act}^{\kappa_e}\). Si l’on compare \(L'=sL\) avec \(\hbar,c,E_e^{(0)}\) fixés, alors \(Q'=Q/s\), \(b_e'=s b_e\) et \(Q'b_e'=E_e^{(0)}\). La masse fixe conserve sa valeur et transporte sa coordonnée lors du changement d’horloge.}{Cela distingue la référence radiale \(L_*\), la base \(u_L=u_Cu_T\) et l’avancée \(\ell_0\). La période \(108t_0\) est \(2\pi L/c\). Le mode \(X=I\) détermine le quantum de l’horloge ; la coordonnée électronique conserve en outre sa normalisation \(b_e=E_e^{(0)}/Q_{\rm ret}\) et son facteur \(R_{\rm act}^{\kappa_e}\). Lorsqu’on compare \(L'=sL\) avec \(\hbar,c,E_e^{(0)}\) fixés, on a \(Q'=Q/s\), \(b_e'=s b_e\) et \(Q'b_e'=E_e^{(0)}\). La masse fixe conserve sa valeur tandis que sa coordonnée est transportée vers la nouvelle horloge.}


\HMTLang{Les deux sections d’action conservent \(G_{\rm pre}/G_{\rm ret}=R_{\rm act}^{-1}\) et \(\ell_P^2=\hbar_aG_a/c^3=L^2\). Dans le système de coordonnées \(L_*=1\,\mathrm m\), \(\mathcal U_S=1\,\mathrm{J\,s}\), avec \(c=299792458\,\mathrm{m\,s^{-1}}\), la composition de retour fournit}{The two action sections retain \(G_{\rm pre}/G_{\rm ret}=R_{\rm act}^{-1}\) and \(\ell_P^2=\hbar_aG_a/c^3=L^2\). In the chart \(L_*=1\,\mathrm m\), \(\mathcal U_S=1\,\mathrm{J\,s}\), with \(c=299792458\,\mathrm{m\,s^{-1}}\), the return composition gives}
\begin{align*}
 L&=1.616254982625126330168262501334561962\ldots\,10^{-35}\,\mathrm m,\\
 t_0&=3.136499962536549073434219319400884017\ldots\,10^{-45}\,\mathrm s,\\
 G_{\rm ret}&=6.674299659414022706367776189182194954\ldots\,
 10^{-11}\,\mathrm{m^3\,kg^{-1}\,s^{-2}}.
\end{align*}
\HMTLang{La section constitutive conserve \(u_C=c/\widehat c\). Ces chiffres évaluent la composition avec les sections déclarées ; les identités normalisées et leurs domaines constituent sa preuve.}{The constitutive section retains \(u_C=c/\widehat c\). These digits evaluate the composition in the declared sections; the normalized identities and their domains constitute its proof.}

\HMTLang{Les échelles positives reconnues dans cette réalisation sont \(t_P=L/c\), \(m_{P,a}=\hbar_a/(cL)\) et \(E_{P,a}=\hbar_a c/L=Q_a\). Substituer~\eqref{xii:radial:G} dans \(\sqrt{\hbar_aG_a/c^5}\), \(\sqrt{\hbar_ac/G_a}\) et \(\sqrt{\hbar_ac^5/G_a}\) démontre respectivement les trois égalités. Pour un mode \(y>0\), \(m(y)=m_{P,a}y\), \(r_g(y)=Ly\) et \(\bar\lambda(y)=L/y\). Leur produit est \(L^2\), tandis que la coïncidence des longueurs caractérise \(y=1\). On a \(t_0=(\pi/54)t_P\) ; un horizon \(R_h=\zeta_hL\) conserve sa variable géométrique \(\zeta_h>0\).}{Les échelles positives reconnues dans cette réalisation sont \(t_P=L/c\), \(m_{P,a}=\hbar_a/(cL)\) et \(E_{P,a}=\hbar_a c/L=Q_a\). Substituer~\eqref{xii:radial:G} dans \(\sqrt{\hbar_aG_a/c^5}\), \(\sqrt{\hbar_ac/G_a}\) et \(\sqrt{\hbar_ac^5/G_a}\) prouve respectivement les trois identités. Pour un mode \(y>0\), \(m(y)=m_{P,a}y\), \(r_g(y)=Ly\) et \(\bar\lambda(y)=L/y\). Leur produit est \(L^2\), tandis que l’égalité des longueurs caractérise \(y=1\). On a \(t_0=(\pi/54)t_P\) ; un horizon \(R_h=\zeta_hL\) conserve sa variable géométrique \(\zeta_h>0\).}

\subsection{\HMTLang{Réduction avec résidu et transport des symétries}{Reduction with a residual and transport of symmetries}}
\HMTLang{L’énergie avec mémoire de la section~\ref{vii:sec:variacion-memoria} conserve le produit hermitien antilinéaire en la première variable, les transports unitaires \(U_a\) et le poids global positif \(W_m\), y compris ses termes croisés. Sa variation complète~\eqref{vii:eq:variacion-completa-memoria} réunit champ, transport \(\dot U_a=A_aU_a\), \(A_a^*=-A_a\), et poids \(\dot W_m\). La preuve par développement au premier ordre et le théorème~\ref{vii:thm:refinamiento-respuesta} restent dans l’article ; ce dernier transporte la réponse au moyen d’inclusions \(J_m,K_m\) fixes. Les variations continues appartiennent au lecteur déclaré ; une route TPK conserve en outre sa réalisation typée. La réduction suivante s’applique à une fibre positive fixée de cette énergie et conserve explicitement son résidu.}{L’énergie de mémoire de la section~\ref{vii:sec:variacion-memoria} conserve le produit hermitien antilinéaire en son premier argument, les transports unitaires \(U_a\) et le poids global positif \(W_m\), y compris ses termes croisés. Sa variation complète~\eqref{vii:eq:variacion-completa-memoria} combine le champ, le transport \(\dot U_a=A_aU_a\), \(A_a^*=-A_a\), et le poids \(\dot W_m\). La preuve par développement au premier ordre et le théorème~\ref{vii:thm:refinamiento-respuesta} restent dans l’article ; ce dernier transporte la réponse au moyen d’inclusions fixes \(J_m,K_m\). Les variations continues appartiennent au lecteur déclaré ; une route TPK conserve en outre sa réalisation typée. La réduction suivante s’applique à une fibre positive fixée de cette énergie et conserve explicitement son résidu.}
\HMTLang{À une coupure commune entre frontière et mémoire, écrivons}{Pour une séparation commune entre frontière et mémoire, écrivons}

\[
 Y=\begin{pmatrix}A&B\\ B^*&C\end{pmatrix}>0,
 \quad S=A-BC^{-1}B^*,\quad r=y+C^{-1}B^*x.
\]
\HMTLang{Compléter le carré et factoriser la matrice démontre}{Completing the square and factoring the matrix proves}
\[
 (x,y)^*Y(x,y)=x^*Sx+r^*Cr,\quad y=r-C^{-1}B^*x,
 \qquad (Y^{-1})_{bb}=S^{-1}.
\]
\HMTLang{Par homogénéité du complément de Schur, les réponses réduites sont}{By homogeneity of the Schur complement, the reduced responses are}
\[
 R_{\rm eff}=LS,\quad H_{\rm eff}=Q_aS,\quad
 \Lambda_b=LS^{-1},\quad R_{\rm eff}\Lambda_b=L^2I,
 \quad H_{\rm eff}\Lambda_b=\hbar_a cI.
\]
\HMTLang{La longueur duale comprime l’inverse ; l’archive \((x,r)\) conserve les états qui partagent un minimum. Une coordonnée inversible de bord \(T\) transporte \(R,H\) par \(T^{-*}(\cdot)T^{-1}\) et \(\Lambda\) par \(T(\cdot)T^*\), en conservant les deux produits. La métrique est transportée conjointement.}{The dual length compresses the inverse; the archive \((x,r)\) retains the states that share a minimum. An invertible boundary coordinate \(T\) transports \(R,H\) by \(T^{-*}(\cdot)T^{-1}\) and \(\Lambda\) by \(T(\cdot)T^*\), preserving both products. The metric is transported jointly.}

\HMTLang{Sous des changements unitaires \(T\) au départ et \(S_0\) à l’arrivée, on transporte simultanément \(B_\Omega'=S_0B_\Omega T^*\), \(P_i'=S_0P_iS_0^*\), \(P'=S_0PS_0^*\), \(X'=TXT^*\) et \(\Delta'(Z)=S_0\Delta(S_0^*ZS_0)S_0^*\). La substitution donne \(D'=S_0DT^*\), \(R'=S_0RS_0^*\), \(H'=S_0HS_0^*\) et \(\Lambda'=S_0\Lambda S_0^*\). Le couplage scalaire demeure fixe. Une représentation \(\rho\) déjà construite sur la fibre d’arrivée se transporte par \(\rho'(g)=S_0\rho(g)S_0^*\), qui conserve toutes ses relations par annulation de \(S_0^*S_0\). Les constructions exceptionnelles précédentes conservent leurs porteurs et leurs applications propres.}{Sous des changements unitaires \(T\) à la source et \(S_0\) à la cible, on transporte conjointement \(B_\Omega'=S_0B_\Omega T^*\), \(P_i'=S_0P_iS_0^*\), \(P'=S_0PS_0^*\), \(X'=TXT^*\) et \(\Delta'(Z)=S_0\Delta(S_0^*ZS_0)S_0^*\). La substitution donne \(D'=S_0DT^*\), \(R'=S_0RS_0^*\), \(H'=S_0HS_0^*\) et \(\Lambda'=S_0\Lambda S_0^*\). Le couplage scalaire demeure fixe. Une représentation \(\rho\) déjà construite sur la fibre cible se transporte par \(\rho'(g)=S_0\rho(g)S_0^*\), en conservant toutes ses relations par annulation de \(S_0^*S_0\). Les constructions exceptionnelles précédentes conservent leurs propres porteurs et applications.}

\HMTLang{Le raffinement conserve une contrainte de phase supplémentaire. Si \(\theta\ne0\) et \(e^{-ia\theta/9^n}=e^{-i\theta/9^n}\) pour tout \(n\), alors \((a-1)\theta/9^n\in2\pi\mathbb Z\) tend vers zéro et est nul à partir d’un certain rang. Par conséquent \(a=1\). L’action relevée et ses phases raffinées conservent la normalisation complète.}{Refinement retains an additional phase constraint. If \(\theta\ne0\) and \(e^{-ia\theta/9^n}=e^{-i\theta/9^n}\) for every \(n\), then \((a-1)\theta/9^n\in2\pi\mathbb Z\) tends to zero and is eventually zero. Thus \(a=1\). The lifted action and its refined phases retain the full normalization.}
